% Discussion moved from Section sec:model of the main paper (sections/model.tex) in revision v2, to shorten the main text.
\section{Further Discussion of the Model}\label{supp:model}
This section collects discussion that accompanies the results of Section~\ref{sec:model} of the main paper.

\subsection{Conventions for growth rates}\label{supp:model-rates}
This subsection accompanies the setup of Section~\ref{sec:model-setup}. The static results can equally be read in period rates, provided the growth technology \eqref{eq:growth} is specified in the convention the score uses; it cannot hold in both at once. If it holds in continuously compounded rates, the frontier with $k=0$ expressed in annual rates turns convex once excess log growth exceeds $\beta-1$ (between 1.2 and 3 for the range of $\alpha$ used in the paper), so a maximal score in annual rates exists only because of the share cap. The two conventions differ materially only at growth rates of the order of 100\% or more: \autoref{prop:saturation}(d) and its Rule-of-1000 example use continuously compounded rates, \autoref{cor:transient} uses period rates, and at scores of several hundred points the binding constraint is share headroom (\autoref{cor:transient}), not the frontier. \autoref{lem:timing} and \autoref{prop:budget}(b) treat the other places where the distinction matters.

\subsection{The two weights on growth in numbers}\label{supp:model-weights}
\emph{The decision weight (\autoref{prop:dynamic}).} Take $r=0.10$, $G=0.04$, $k=0$, $\alpha=0.3$ and $m_0=0.35$, so that a firm that never spends on growth is worth $m_0/(r-G)\approx 5.8$ times revenue. Demand pull $\lambda\approx 0.070$ makes the small-share multiple $v_0=10$: the firm grows at 7.6\%, spends 10.7\% of revenue on growth and earns an FCF margin of 24.3\%, and $0.243/(0.10-0.076)=10$. Demand pull $\lambda\approx 0.081$ gives $v_0=15$, with growth of 9.2\%, spend of 23.6\% and a margin of 11.4\%. Holding $\lambda\approx0.070$ and changing only the discount rate, the small-share multiple, the smaller root of $(r-G)v_0=m_0+\frac{1-\alpha}{\alpha}(\alpha\lambda v_0)^{1/(1-\alpha)}$, is \WeightWacc, \WeightTen, \WeightTwelve, \WeightFifteen, \WeightTwentyTwo{} and \WeightThirtyOne{} at $r=9.34$, 10, 12, 15, 22 and 31\%: an owner who discounts at 22\% or more acts on a weight within or below the screening range of 2 to 3. With the share cap $\bar s=0.5$ and $v_0=10$, a numerical solution of \eqref{eq:hjb} gives a growth weight $v'(s)$ of 9.1, 7.2 and 5.4 at 1\%, 10\% and 50\% of the cap, against revenue multiples $v(s)/s$ of 9.4, 8.1 and 6.5. In this example the firm's own weight is close to, and below, its revenue multiple, and it falls with share. It exceeds the screening weights of 2 to 3 when the firm's revenue multiple does; a firm valued near three times revenue would have a decision weight near or below three. Equity cash flows of longer duration earn lower expected returns \citep{weber2018cash,gormsen2023duration}; if the discount rate falls with duration, the marginal $q$ of revenue of a fast grower is higher than this constant-rate calculation gives.

\emph{The screening weight (\autoref{lem:valuation2}).} The screening weight is a duration, as in the implied equity duration of \citet{dechow2004implied}: it rises with how long growth persists and with how much of the value lies in the future, and it falls with the persistence of the current margin. For example, $m=0.25$, $D=10$ and $g=0.25$ give $X^{\mathrm{scr}}=0.25\cdot 10/1.25=2$, inside Bessemer's range; with $r=0.10$, a duration of 10 years at $g=0.25$ requires about 15 years of 25\% growth ($D=10.25$ at $T=15$, but only $D=6.5$ at $T=10$). Part (b) makes $X^{\mathrm{scr}}$ proportional to the current margin, so a firm that breaks even because it spends heavily on growth would get no credit for growth; part (c) removes this artifact. With $r=0.10$, $g=0.25$, $G_\infty=0.03$, $\bar m=0.25$, $T=10$ and $\theta=1/2$, \eqref{eq:X2} gives $X^{\mathrm{scr}}\approx 1.8$ for $m=0.10$ and $2.8$ for $m=0.25$; at $\theta=0$ these rise to 5.4 and 6.2, and with $T=1$ and $\theta=0$ the weight for $m=0.10$ is 3.0, the revenue multiple times $(1+r)/(1+g)^2$.

\emph{Equal score, unequal value.} Table~\ref{tab:isoscore} applies \autoref{lem:valuation2} to four firms on the same Rule-of-40 line. At $r=9.34\%$, ten years of growth and the current margin held throughout, a firm growing 30\% at a 10\% margin is worth \IsoThirty{} times revenue and one with no growth and a 40\% margin \IsoZero{} times. The table is an implication of the lemma at stated inputs; the paper estimates no mapping from scores to market multiples.

\IfFileExists{sections/tables/tab_isoscore.tex}{\begin{table}[t]\centering\small
\caption{Four firms on the same Rule-of-40 line are worth different revenue multiples (Lemma~\ref{lem:valuation2}, an analytical illustration)}\label{tab:isoscore}
\begin{tabular}{@{}rrrr@{}}\toprule
Growth $g$ (\%) & Margin $m$ (\%) & Rule-of-40 score & Value $V/S_0$ (times revenue)\\\midrule
30 & 10 & 40 & 12.1\\
20 & 20 & 40 & 11.7\\
10 & 30 & 40 & 8.3\\
0 & 40 & 40 & 5.2\\
\bottomrule\end{tabular}
\par\vspace{2pt}\footnotesize\raggedright Notes: $V/S_0=m\,\mathcal N_T(\eta)+m\,\eta^T\mathcal M_\infty$ with $\eta=(1+g)/(1+r)$, growth $g$ for $T=10$ years and 3\% thereafter, the current margin held throughout ($\theta=1$), and $r=9.34\%$. Every row scores 40. The table is an implication of the lemma at stated inputs; the paper estimates no mapping from scores to market multiples.
\end{table}}{}

\subsection{Reading the cross-sectional slope under a common frontier}\label{supp:model-slope}
This subsection accompanies the opening of Section~\ref{sec:model-industry}. If every firm has the frontier $m^*$ and firms act on different weights, their chosen points lie on $m^*$, and $\Psi_\chi(g)$ is $m^*$ evaluated at the $(1-\chi)$-quantile of growth among firms growing at least $g$, so its slope approximates that of $m^*$ when growth rates are dense. With $k=0$ and the growth window read as excess growth $\Delta\in[0,W]$, the frontier $m_0-(\Delta/\ell)^{\beta}$ has least-squares slope $-b$ on an equally spaced grid of the window if and only if $\ell=[S_\beta/b]^{\alpha}$, where $S_\beta$ is the least-squares slope of $\Delta^{\beta}$ on that grid, approximately $6\beta W^{\beta-1}/[(\beta+1)(\beta+2)]$; its steepest slope in the window, at $\Delta=W$, is then about $(\beta+1)(\beta+2)\,b/6$, between $2.3b$ and $5b$ for $\alpha\in[0.25,0.45]$.

\emph{An example of selection by quality.} A low type with $(\ell,m_0)=(0.3,0.30)$ and a high type with $(0.8,0.50)$, with $\alpha=0.3$, $X=1$, $k=0$ and $G=0.1$, choose $(g,m)\approx(0.21,0.27)$ and $(0.53,0.37)$: the cross-sectional slope is about $+0.31$, while both firms' frontiers have slope exactly $-1$ at their chosen points.

\subsection{The assumptions behind free entry}\label{supp:model-entry}
This subsection accompanies \autoref{prop:entry}. The result is specific to its assumptions, three of which do most of the work. Entry happens only at date 0: with $G>0$ flow profits grow without bound, so later entry would be profitable for any $F$, and allowing it would drive the markup to zero faster, not slower; the conclusion is conservative in that direction only. The scarce input's supply grows at exactly the market rate, a knife-edge case that keeps its price constant; with slower supply growth the rent would rise over time. And the product is homogeneous: with differentiated products the markup stays at the level that substitutability sets (\autoref{rem:bertrand}). Nor does the model let firms respond to cheaper replication by spending more on quality or growth, which can keep concentration high when such sunk outlays are chosen, as in the theory of endogenous sunk costs \citep{sutton1991sunk}. Rents can also accrue upstream: the supply of foundation models has strong economies of scale and scope \citep{korinek2025concentrating}, so a scarce input may belong to the model provider and appear as a markup inside $c_c$ ($\upsilon_c$ in \autoref{rem:owncompute}) rather than as the software firm's $\rho$.

\subsection{What each cost floor contains}\label{supp:model-floors}
This subsection accompanies \autoref{def:fullcost}. The companion paper \citep{murthy2026transformation} derives the same floors activity by activity. Each cost is counted once. The accountability floor $a_{\min}$ keeps its meaning (officer certification, internal control over financial reporting and its integrated audit, directors' and officers' insurance, legal and liability costs, and model-risk validation, which here means the firm's own validation: the U.S. banking agencies' revised guidance of April 2026, which superseded SR 11-7, places generative and agentic AI models outside its scope; \citealp{fed2026sr262,occ2026bulletin13}), and it also holds whatever human review and sign-off regulation or the firm's risk appetite requires. By \autoref{prop:amin} it further holds the control layers and IT general controls that a control design meeting the material-weakness standard requires once AI performs financial controls, and the expected cost of errors below materiality. The same model upgrade triggers two different activities, each counted once: the functional revalidation of every automated build (re-running the evaluations and regression tests that show the build still performs its task) is in $c_{\mathrm{mnt}}$; the retest of the IT general controls of an in-scope application (evidence that the change-management, access and application controls over the upgraded model still operate) is in $a_{\min}$, and covers only the control testing beyond the functional revalidation. Operational losses that survive review are in $c_{\mathrm{err}}$, and neither it nor its common-mode term counts misstatements; incidents caused by attackers, including injected instructions and exploited shared vulnerabilities, are in $c_{\mathrm{sec}}$. One-time costs (builds, parallel running, severance, training) are transition costs and enter Section~\ref{sec:model-attainable}, not the ceiling.

\subsection{Common-cause failure across providers}\label{supp:model-commonmode}
This subsection accompanies \autoref{ass:layers} and \autoref{prop:amin}. The assumption treats failures of different types as independent apart from the cross-type term $p^{\mathrm{cc}}$. That is a simplification, and it favors diversity. Errors of models from different providers are correlated: across 71 models on 14,042 four-option MMLU questions from HELM, when two models are both wrong they give the same wrong answer about 60\% of the time on average, against one third for uniformly random wrong answers \citep{kim2025correlated}. That figure is the probability of the same wrong answer given that both models are wrong; it is not the probability that one model errs given that the other does, which it does not measure. For a control that flags an item when two providers disagree, it means that about 60\% of joint errors pass unflagged. We therefore read the value of diversity in \autoref{prop:amin}(a) as an upper bound, and we treat a second provider as a separate type only to the extent that its failures are not common with the first. The common-mode probabilities are also conditional on the IT general controls: a firm that weakened those controls would raise them. The cross-type term covers management override, defective shared data, failures of the IT general controls and reviewers' acceptance of wrong AI output, not shared tooling as such: an identity provider, an endpoint security agent or a cloud region that every layer depends on can defeat layers from different providers at once. The model has no term for that dependence, so the diversity rule addresses the failures of models, not the availability and tooling risk of the infrastructure beneath them, and the full-cost floor does not include the cost of keeping a second provider ready.

\subsection{The program clock}\label{supp:model-clock}
This subsection accompanies \autoref{def:attainable}. Three features of part (b) are assumptions, and each has a stated reason. The uniform clock is a mean-field reading of reorganizations: \autoref{prop:expected}(f) shows that it is the almost-sure long-run limit of a clock with random freezes, and states the results that hold under random freezes at every date. Disruptions run in calendar time because the turnover and strategy data that ground $\nu_D$ are rates per calendar year. And a realized share that a second migration does not change describes the shortfall of a migrated process as a property of the process (how much of the targeted value the automated work actually delivers), which is how the calibration's source measures it: the share of targeted value that transformations had realized once they were done \citep{mckinsey2021losing}. Disruptions are independent across attempts; if a disruption instead ends every attempt in flight at once, the bounds and equalities of \autoref{prop:expected}(c) and (e)(i)--(ii) are unchanged, because they depend only on per-attempt probabilities, but the variance of the embedded share and of the dip is higher.

Part (e) is the version that the calibration of Section~\ref{sec:horizons} computes, with $c_w=\rho T_{\mathrm{par}}$ for a binding capacity $\rho$ expressed as a share of the program a year: the expected realized share $\bar b\min\{1,p^{\rm emb}\rho T_{\mathrm{par}}n(t_R)\}$ is an upper bound that the greedy plan attains while every wave can be filled, and overstates by less than $\bar b\rho T_{\mathrm{par}}$ afterwards. Part (f) shows that the uniform clock is a long-run approximation: at a given date, random freezes give a dispersed program time, and the expected number of embedded waves is not the number at the expected program time. \autoref{lem:expected-frontier} in the Supplementary Appendix adds that, when processes shift only $m_0$ and $k$, the expected frontier is the frontier at the expected parameters, and that it understates the expected maximal score by the value of choosing growth after outcomes are known.

\subsection{The compute share, its path and pass-through}\label{supp:model-compute}
This subsection accompanies Section~\ref{sec:horizons-ceiling}. \emph{Evidence on the compute share.} On METR's software tasks more than 80\% of successful runs cost less than 10\% of the human cost at \$143.61 an hour \citep{kwa2025measuring}, and that runs cost 0.4 to 1.1\% of it, rounded to 1\%, is our computation from METR's public run data \citep{metr2026evalanalysis}; this sets the lower credible bound of \CcLo\%. A survey puts the cost of revenue of AI products sold to customers at 55\% in 2025 \citep{iconiq2026state}, but that includes talent and hosting and so overstates the compute cost of a firm's own automation; the same survey puts spending on AI for internal productivity, direct and indirect and not itemized, at 11\% of revenue in 2025.

\emph{Drivers of the path.} Two scenario values drive the path of the compute share: the price of the work that only frontier models reach near each horizon, \CcKappaFront{} of the labor cost it replaces, within the value test, and the share of automated work that only frontier models reach, \CcOmegaTwentyEight, \CcOmegaThirtyOne{} and \CcOmegaThirtySix\% at the three horizons. The frontier-priced share falls as the feasible set grows; after 2028 the measured price declines hardly bind, so the compute share is about a fifth of that share. With the frontier share in 2028 at 30\% or 90\% the 2028 threshold is \GstarOmegaLoTwentyEight{} or \GstarOmegaHiTwentyEight\%.

\emph{Pass-through.} By \autoref{lem:prices}(a) the elasticity $\epsilon_p$ that corresponds to an industry pass-through rate $\varphi$ is $\varphi u/p$, \PtElasHalf{} at $\varphi=0.5$. By horizon, with $\varphi=0.5$, the thresholds are \GstarPhiHalfTwentyEight, \GstarPhiHalfThirtyOne{} and \GstarPhiHalfThirtySix\%. A firm keeps the share $1-\varphi(1-H)$ of a common dollar saving once its gain in volume is counted, at a concentration $H$ \citep{murthy2026transformation}, but with costs per unit its margin does not depend on volume, so $H$ does not enter the score. Table~\ref{tab:passthrough-recon} gives the threshold under each reading.

\subsection{Execution: yield, ownership and the change from the first version}\label{supp:model-execution}
This subsection accompanies the pace bullets of Section~\ref{sec:horizons-ranges}. In \autoref{prop:expected} only embedded transitions realize value, so the targeted savings that an attempt realizes in expectation, its yield, is $\varpi\bar b$: \YieldModel{} under central execution. The share $\varpi=\PaceSuccCen$ rounds down the share of transformations in one survey that created value, the win zone (targets met) and the worry zone (some value, targets missed) together \citep{bcg2020flipping}; $\bar b=\PaceRealCen\%$ is a second survey's realized share for successful transformations, against 37\% for all others \citep{mckinsey2021losing}. Combining the first survey's win-zone share of 30\% with the second survey's shares, so that every transformation outside the win zone realizes 37\%, gives an unconditional yield of \YieldMck{} and a realized share by 2036 of \ShareMck\%, close to the central path. Read like for like, with the worry zone realizing 37\% and the woe zone, which created no value, realizing nothing, the yield is \YieldExec{} and the realized share by 2036 \ShareExec\%. If only the win zone counts, at the realized share of a success, the realized share by 2036 is \ShareBcgWin\% and the 2036 score \ScoreBcgWin, a row of Table~\ref{tab:tornado}. The surveys cover whole transformations, not process-level AI attempts, so all three readings are analogies.

Under the execution inputs of the first version of this paper the central firm would realize \BridgeOld\% of its program's benefits by 2036, against \BridgeNew\% now. A Shapley decomposition over the three changes, which sums exactly, attributes \BridgeSucc{} points to a higher share of attempts that create value and the removal of a separate disruption haircut (\BridgeVarpi{} and \BridgeHaircut{} respectively), \BridgeBbar{} to a lower realized share of a success and \BridgeTemb{} to a shorter embedding time.

\emph{Program structure.} Change capacity maps one for one onto the program: changing the work of \PaceCapCen\% of employees a year covers \PaceCapCen\% of the program a year, which assumes that the program spans all employees' work evenly, as full automation approximately does. The calibration treats the program as a continuum of identical processes (\autoref{prop:expected}(e)), so its realized share rises in proportion to the capacity spent. When savings are concentrated in a few processes, a plan that attempts first the processes with the largest gain per unit of the binding capacity, as the bound of \autoref{prop:attainable}(d) orders them, realizes more of the program's benefits early, and the central path understates it. That change capacity binds follows from the scenario values for the program's size, decision bandwidth and security review; engineering and data capacity are not modeled as resources, and if either binds the path is slower.

\emph{Ownership and slow execution.} A public firm that holds to a guided dip of one point of revenue a year can attempt only about \VarDipOneInit\% of the central program at first, and its budget loosens only as savings accumulate; by 2036 its intercept rises by \VarDipOneMzero{} points instead of \VarCenMzero. Faster decisions matter only for programs whose decisions are slow relative to the number of employees affected. An owner can also raise change capacity, for example by replacing staff, or accept a larger dip than a public firm guides to; the central case models neither, and this paper does not analyze private-equity or permanent-capital owners. Under slow execution (\PaceCapSlow\% of employees a year, a share 0.54 of attempts creating value, 37\% of the targeted savings realized, frequent disruption), \VarSlowShare\% of the program's benefits are realized by 2036.

\subsection{Scores in filings}\label{supp:model-filings}
This subsection accompanies Section~\ref{sec:discussion}, ``Use in filings''. A score computed on adjusted margins, on free cash flow, on organic growth or with acquisitions charged as capital is a non-GAAP financial measure when an issuer makes it public. It then needs the most directly comparable GAAP measure, presented with equal or greater prominence, and a quantitative reconciliation \citep{cfr2026sec244100,cfr2026sec22910,sec2022cdi}. A performance measure may not remove items as non-recurring that are reasonably likely to recur within two years, and a measure that excludes normal, recurring cash operating expenses can be misleading; for a serial acquirer that question extends to acquisition and reorganization costs. Free cash flow is a liquidity measure and may not be presented per share \citep{sec2022cdi}. Palantir defines its Rule-of-40 score under ``Additional Definitions'', outside its non-GAAP section, and reconciles the adjusted margin it uses but not the score \citep{palantir2026q2}. A Rule-of-$X$ metric discussed in management's discussion and analysis is also a key performance indicator, for which the Commission expects a clear definition and calculation, an explanation of why it is useful and how management uses it, the estimates and assumptions behind it, disclosure of changes in its method and effective disclosure controls \citep{sec2020release3310751}. The ceiling, the durable score, the growth surplus, the decision weight and the central path rest on unobserved parameters and scenario values and are not suitable as reported metrics or targets; stated as a company's expectations they would be forward-looking statements, for which the statutory safe harbor is not available in an initial public offering \citep{usc2026sec77z2}.

\subsection{Ownership regimes}\label{supp:model-regimes}
\begin{remark}[Ownership regimes]\label{rem:regimes}
An owner sets the weight $X$, the discount rate, the horizon and the tolerated dip, and its decision speed scales the decision capacity; the companion paper \citep{murthy2026transformation} derives these for public, private, private-equity and other owners. On a common technology, three consequences follow. (i) Owners with different weights choose different points of the same attainable set: a higher weight picks a point with weakly more growth, and, when the lower weight is positive and growth is strictly higher, less margin. (ii) An owner who admits only the transitions that pay within its rate and horizon admits a subset $\mathcal P'\subseteq\mathcal P$, whose attainable and static sets are contained in those of an owner who admits all of $\mathcal P$ with the same capacities. (iii) Pace differs: a tighter dip budget, or slower decisions that lower decision capacity, weakly shrink $\mathcal A_t$ and weakly lower $\bar R_X(t)$ at every horizon (Propositions~\ref{prop:attainable}(a) and \ref{prop:expected}(b)), while the static frontier of a given program is unchanged.
\end{remark}
